\makeabstract
    {
    Long Story Short: Metal identification that transfers from short peptides to long
    }
    {
    Sarah Kong\textsuperscript{1}, Jasper Steenwinkel\textsuperscript{2}, Jeffrey Vargas\textsuperscript{3}, Sunera Wijeratne\textsuperscript{4}, Theresa L. Windus\textsuperscript{4}
    }
    {
    \textsuperscript{1} Department of Chemistry, Amherst College, Amherst, MA\\
\textsuperscript{2} Department of Chemistry, Ripon College, Ripon, WI\\
\textsuperscript{3} Department of Chemistry, Arizona State University, Tempe, AZ\\
\textsuperscript{4} Department of Chemistry, Iowa State University, Ames, IA
    }
    {
    Transition metals are underutilized in electronic and industrial waste, and peptides offer a practical way to reclaim them because they can be tuned to bind and separate specific metal ions. Identifying which metal a given peptide has bound is necessary for this application, and infrared (IR) spectroscopy offers an inexpensive way to do so, since each complex produces a distinct vibrational fingerprint. Because longer peptides are more expensive to simulate and measure, this study tests whether a model trained only on short metal-peptide complexes (monomers, single amino acids; and dimers, two-residue peptides) can generalize to longer, unseen complexes (trimers, three-residue peptides).
    }
    {
    IR spectra were generated in bulk using semiempirical quantum chemistry (xTB), avoiding costly experimental synthesis and measurement before validating candidates. A total of 128,854 spectra were simulated for cobalt, nickel, copper, and zinc complexes across monomer, dimer, and trimer peptide lengths (24,685 mono/dimer; 104,169 trimer), each represented as 791 binned intensity values (50{\textendash}4000 cm\ensuremath{^-}{\textonesuperior}). A random forest classifier was trained on mono/dimer spectra only and evaluated via repeated Monte Carlo cross-validation (10 seeds) using macro-averaged F1 score (macro-F1) as the primary metric. The trained model was then tested zero-shot on held-out trimer spectra, including a size-matched trimer subset to control for test-set-size effects. Follow-up experiments varied the composition and volume of trimer data added to training.
    }
    {
    The model classified metal identity at 0.81 macro-F1 on held-out mono/dimer spectra (chance = 0.25), with nickel and cobalt the most confused pair, consistent with their similar binding geometry and molar mass. Zero-shot transfer to trimers dropped performance to 0.63 macro-F1, a drop confirmed to reflect the length shift rather than test-set size alone. Adding a modest, diverse set of distinct trimers to training closed most of this gap, whereas adding duplicated trimer data or additional short-peptide data did not; a model anchored on short peptides plus added trimers reached the same performance plateau ({\textasciitilde}0.73 macro-F1) as one trained directly on trimers.
    }
    {
    Metal identity encoded in computed IR spectra of short metal-peptide complexes substantially transfers to unseen, longer complexes, and the transfer gap closes more efficiently by diversifying training data across peptide length than by increasing data volume alone. This points to a cost-efficient strategy for extending machine-learning-based metal identification to longer peptides: anchor training on abundant short peptides and supplement with a modest, diverse set of longer examples.
    }

    \newpage
    